% Merged into the Experimental Analysis and Discussion chapter.

\section{Interpretation of Results}

The results show that a food recognition system is useful only when detection is connected to the rest of the workflow. The detector tells the system what appears in the image, but the nutrition layer decides what that label means in calories and macros. The frontend then makes those numbers visible and adjustable. The main outcome of this project is therefore the integration of these layers, not only the training of a YOLO model.

The dataset engineering work was also important. A model trained on inconsistent labels would produce confusing outputs even if training completed successfully. Canonical class normalization gave the detector a stable label space and also made nutrition mapping possible. The visual verification sheets were useful because they exposed dataset issues that a file-count check would miss. A few noisy annotations remain, but the class list itself is aligned.

One practical observation from the project is that food recognition errors are not always clear from metrics alone. A small class-label error becomes more serious when it leads to a wrong nutrition record. For this reason, the results have to be interpreted as a complete chain from image to recommendation, not only as detector scores.

\section{Strengths of the System}

One strength of the system is its full-stack completeness. The project does not stop after model inference. It includes dataset building, backend APIs, nutrition database integration, recommendation logic, and a frontend workflow that can be used and inspected. This made the implementation more difficult, but it also made the final output closer to a real application.

Another strength is domain specificity. The dataset focuses on Indian foods, and the nutrition database uses Indian food names and macro values. The alias and mapping logic also accounts for regional naming patterns, which matters for labels such as \code{satham}, \code{medu_vadai}, \code{thengai_chutney}, and \code{nandu_masala}.

The conservative mapping strategy is also a strength. Classes are mapped through explicit semantic rules first, and classes absent from INDB are covered only when a source-backed supplemental row is available. In a nutrition-related application, this reduces the risk of weak fuzzy matches producing misleading macro values.

The project is also documented through reproducible artifacts. It contains scripts for dataset construction, nutrition database creation, visual class verification, and macro verification. These scripts make the data pipeline inspectable instead of leaving the dataset and mappings as unexplained files. For a final-year project, this is useful because the work can be checked through generated artifacts rather than only through screenshots.

\section{Limitations}

The main limitation is portion size. The system reports macros per 100 grams and allows serving adjustment, but it does not estimate exact food mass from the image. This limits calorie precision. Accurate portion estimation would require depth information, reference objects, segmentation masks, or learned volume estimation.

Nutrition source consistency is another limitation. Five classes are covered through supplemental external sources because the available INDB sheet lacks safe equivalent entries. Some other classes use documented substitutes rather than exact dish rows. This is acceptable for a prototype, but a health-sensitive deployment would need stronger food-composition validation and expert review.

Dataset noise also remains. Visual verification found a small number of annotation issues. Large food datasets often contain such noise, but object detectors are sensitive to label quality. In this project, noisy samples did not prevent a working detector from being trained, but they should still be considered while interpreting the results.

The detector also behaves unevenly across visually similar dishes. The final YOLO11s metrics are suitable for a working prototype, but the confusion matrix shows that some classes are harder than others. This is especially important for a nutrition application because a wrong food class can propagate into an incorrect nutrition lookup even when the interface and database are functioning correctly.

The recommendation layer depends on an external language model service when configured. This is acceptable for a prototype, but it makes recommendation output harder to validate than the rule-based macro calculation. A self-trained recommendation model with fixed input-output constraints would provide better control in a later version.

\section{Ethical and Practical Considerations}

Nutrition applications can influence user behaviour. Therefore, the system should avoid presenting estimates as medical advice. Recommendation text should be framed as general guidance, and users with medical conditions should consult qualified professionals. The application should also communicate uncertainty when a detected food uses supplemental nutrition data or when model confidence is low.

Data privacy is another consideration. Food images can reveal personal habits, location context, and health patterns. A production version should minimize image retention, protect API communication, and provide clear user consent. The current prototype is local and does not implement cloud storage for uploaded images.

There is also a practical responsibility in how results are displayed. A bounding box and macro card can make an output look certain even when the food class or serving size is uncertain. For this reason, confidence scores, nutrition sources, and serving controls are important parts of the interface.

\section{Why Supplemental Rows Were Added}

The classes \code{bhakarwadi}, \code{ghevar}, \code{jalebi}, \code{khandvi}, and \code{nandu_masala} were absent from the available INDB file. Mapping them through fuzzy matching would be unsafe because a fried sweet should not be mapped to a generic sweet if the macro profile is different, and a crab masala class should not be mapped to a spice mix or unrelated seafood item. The database therefore adds verified supplemental rows for these classes and records their source URLs. This gives complete dataset coverage while keeping provenance visible.

This decision follows a rule used throughout the project: it is better to document an approximate or supplemental value than to hide a weak match behind a confident interface. In a nutrition application, transparency is more important than making every lookup appear equally reliable.

\section{Lessons Learned}

The project shows that applied AI systems require more than model selection. Dataset curation, naming conventions, data contracts, API design, error handling, and UI workflow affected the final system as much as the detector choice. Food recognition is especially dependent on domain knowledge because visual labels and nutrition records rarely use identical names.

Another lesson is that validation artifacts are valuable. The class verification sheets made it possible to check 72 classes quickly and gave a simple way to discuss dataset quality without opening hundreds of files manually. Without such artifacts, label errors could remain hidden until model performance drops.

The most practical lesson was that every stage needs an independent way to be checked. The dataset script, nutrition validation endpoint, macro reference tables, and frontend workflow checks were not extra polish; they were necessary for finding where a wrong output came from. This made the final report more evidence-based and made the system easier to explain.
